\subsection{A positive-rank primitive-zero cover without a supersingular factor}

\begin{claim}{B.LIT7}{Proved}{An exact cyclic-thirteen distinction between B4 and B5}
Let $C/\F_2$ be the smooth proper model
\[
 C:\quad Y^2+X^2Y=X^5+X,
 \qquad P=\infty.
\]
It has genus two, $2$-rank one, and no nonzero supersingular
isogeny factor in its geometric Jacobian. Over $k=\F_{64}$
there is a connected $P$-pointed cyclic \'etale cover
$D\to C_k$ of degree $13$, with zeta numerator
\[
 P_D(T)=
 (1-T-80T^2-64T^3+4096T^4)
 (1+12T+64T^2)^{12}.
 \tag{B.LIT.7}
\]
Its primitive degree-$13$ factor has Newton slopes
$1/3$ and $2/3$, each with multiplicity $12$, and hence
has $2$-rank zero. Its relative class number is
\[
 \frac{|J_D(k)|}{|J_C(k)|}=77^{12}
       =43439888521963583647921,
\]
which is odd. Nevertheless $D$ has no nonconstant map,
even over $\overline{\F}_2$, to any supersingular curve
of positive genus. Thus this is a B1 witness and a
nontrivial positive-$2$-rank case of B4, but this particular
cover does not supply the correspondence asked for in B5.
\end{claim}

\begin{proof}
The substitution $x=X^{-1}$, $y=Y/X^2$ identifies the
function field with
\[
 y^2+y=x^3+x^{-1}.
\]
The Artin--Schreier poles have odd orders three and one.
Their different exponents are four and two, so
Riemann--Hurwitz gives genus two. In the polynomial model
the regular differential $dX/X^2$ has divisor $2P$.
Direct counts give $|C(\F_2)|=4$ and $|C(\F_4)|=4$;
the genus-two functional equation therefore gives the
Frobenius characteristic polynomial
\[
 H_2(U)=U^4+U^3+2U+4.
\]
It is irreducible modulo three and hence over $\Q$.
Its Newton slopes are $0,1/2,1/2,1$, which also verifies
the asserted $2$-rank. Taking sixth powers of its roots
gives
\[
 H_{64}(U)=U^4-U^3-80U^2-64U+4096,
 \qquad |J_C(k)|=H_{64}(1)=3952=16\cdot13\cdot19.
\]

Here is a complete definition of the cover used in the
computation. Write $k=\F_2[z]/(z^6+z+1)$ and encode
elements by the binary coefficients in the basis
$1,z,\ldots,z^5$. In reduced Mumford coordinates for the
displayed polynomial model, put
\[
 G=\bigl(X^2+25X+14,\;17X+26\bigr)\in J_C(k).
\]
The numbers in this formula are binary field encodings,
not integers embedded in a characteristic-two field.
The certificate verifies $G\ne0$ and $13G=0$.
Since the $13$-primary subgroup has order $13$, there is
a unique surjection
\[
 \chi:J_C(k)\longrightarrow\Z/13,
 \qquad [304]A=\chi(A)G.
\]
Pull back the Lang isogeny $F_{64}-1:J_C\to J_C$
along $R\mapsto[R-P]$ and quotient its deck group
$J_C(k)$ by $\ker\chi$. This defines $D$. The point
above $P$ represented by zero makes it pointed. The
Abel map induces the isomorphism from the abelianized
geometric fundamental group of $C$ to the fundamental
group of $J_C$; thus the pullback of this connected
cyclic isogeny cover is geometrically connected.
This is the precise Jacobian covering statement of
Milne, \S9, Proposition~9.1, applied after base change
to $\overline{k}$: that proposition assumes a separably
closed field, a smooth proper curve, and a chosen point.
\citep[\S9, Proposition~9.1]{MilneJac}

For every $R\in C(k)$ the program computes
$\chi([R-P])$ by exact Cantor arithmetic. Its complete
table, including every point and every element of
$\langle G\rangle$, gives
\[
 \#\{R:\chi([R-P])=0\}=16,
 \qquad
 \#\{R:\chi([R-P])=j\}=4
 \quad(1\le j\le12).
\]
Every polynomial division in the group law is checked
to be exact, and every output Mumford pair is checked
against the curve equation. The producer
\path{code/b5_char2_13_certificate.py} writes these finite
tables and the entire expanded degree-$28$ numerator to
\path{code/b5_char2_13_certificate.json}; its textual
output is also retained.

The independent verifier
\path{code/c0_cyclic13_verify.py} imports no Cantor
implementation. It uses fixed-point addition by the
Chinese remainder theorem and tangent Taylor expansion,
and independently reproduces all $64$ point projections,
all $13$ subgroup elements, and $H_{64}$ by a Newton-sum
recurrence. Its exact output and certificate are retained
beside the producer.

Let $\zeta$ be a primitive thirteenth root of unity.
The complete twisted polynomial has degree two. Its
linear coefficient, for every $i\ne0$, is
\[
 \sum_{R\in C(k)}\zeta^{i\chi([R-P])}
       =16+4\sum_{j=1}^{12}\zeta^{ij}=12.
\]
Because the canonical divisor is $2P$, its quadratic
coefficient is $64$. These are exactly the effective-divisor
and functional-equation formulas proved in B.CL2, so no
single evaluation at $T=1$ is substituted for a Frobenius
polynomial. Multiplying the twelve complete twisted
polynomials and the numerator of $C_k$ proves (B.LIT.7).
It also gives the relative class number and
$g(D)=13(g(C)-1)+1=14$.

The two roots of $U^2+12U+64$ have $2$-adic valuations
two and four: the point $(1,2)$ lies strictly below the
segment joining $(0,6)$ and $(2,0)$. Dividing valuations
by six gives slopes $1/3$ and $2/3$. Thus the primitive
factor has $2$-rank zero and no supersingular factor.
For completeness, $J_C$ has no geometric supersingular
factor either. If a Frobenius root $\alpha$ belonged to
one after a finite constant extension, then
$\alpha/\sqrt2$ would be a root of unity. One elementary
way to see this criterion is that a supersingular Weil
number divided by the appropriate square root of the
field size is an algebraic unit at every finite place
and has absolute value one at every complex place;
Kronecker's argument then makes it a root of unity.
The property descends from $\alpha^m$ to $\alpha$.
Irreducibility of $H_2$ would make all four of its roots
have this property, contradicting its slope-zero root.
This also proves geometric simplicity here: a geometric
elliptic decomposition of a $2$-rank-one surface would
include one supersingular elliptic factor.

Finally $J_D$ is isogenous to $J_C\times\Pr(D/C)$
by pullback and norm. A nonconstant map $D\to S$ of
smooth proper curves induces pullback and norm with
composition multiplication by its degree on $J_S$.
Consequently a positive-dimensional supersingular
$J_S$ would give a supersingular factor of $J_D$,
which the preceding two paragraphs exclude. This holds
for inseparable as well as separable maps, since a
nonzero multiplication map is invertible in the
isogeny category.
\end{proof}

\paragraph{The scope of the search.}
This does not refute the existence of a different odd
\'etale cover of $C$ with a supersingular factor. The
independent theta-Taylor elimination in
\path{code/b5_theta_division.py} excludes torsion orders
$3,5,7,9,11$ on the Raynaud theta curve of this example;
the order-$13$ candidate locus is $r^2+r+1=0$.
The companion probe records complete character sums
for exponent-one prime quotients over all
$\F_{2^m}$ with $1\le m\le10$. These bounded searches
are retained as evidence about the attempted route,
not as a classification of all odd covers or all
constant extensions.

\subsection{Two exact limits of the proposed correspondence constructions}

\begin{claim}{B.LIT8}{Proved}{Ordinary cubic Pryms in two small genera}
Let $C$ be a smooth hyperelliptic curve over an algebraically
closed field of characteristic two, of genus $g\in\{2,3\}$
and $2$-rank zero. For every connected cyclic \'etale
cover $D\to C$ of degree three, the Prym variety is
ordinary of dimension $2(g-1)$. In particular,
$f_D=2(g-1)$. In genus three its genus-two reflection
quotient is ordinary, so this cubic-reflection construction
does not yield a supersingular target.
\end{claim}

\begin{proof}
The hyperelliptic Artin--Schreier map has exactly one
branch point, since its number of branch points is
$f_C+1=1$. Put this point at infinity, denote its
unique preimage by $P$, and use an equation
$y^2+y=f(x)$ with $\deg f=2g+1$. Then
\[
 \operatorname{div}(dx)=(2g-2)P,
 \qquad B^1\simeq\mathcal O_C((g-1)P),
\]
where $B^1$ is Raynaud's line bundle of locally exact
differentials on the Frobenius twist. We identify that
twist with its corresponding equation over the perfect
ground field; all arguments below are unchanged on it.

Suppose a nonzero three-torsion class $L$ satisfies
$H^0(B^1\otimes L)\ne0$. Removing copies of $P$ and
hyperelliptic conjugate pairs from an effective divisor
of degree $g-1$ gives
\[
 L\simeq\mathcal O_C(E-dP),\qquad 1\le d\le g-1,
\]
with $E$ reduced in the hyperelliptic divisor sense:
it contains no $P$ and no conjugate pair. There is a
rational function $h$ with
$\operatorname{div}(h)=3E-3dP$. For $g=2$ or $3$,
\[
 3d\le3(g-1)\le2g<2g+1.
\]
The pole orders of $x,y$ at $P$ are $2,2g+1$.
Consequently every function whose only pole is at
$P$ of order at most $2g$ belongs to $k[x]$.
Every finite zero of such a function occurs together
with its hyperelliptic conjugate and with the same
multiplicity; there are no finite branch points.
This contradicts the form of $E$. Thus neither $L$
nor $L^2$ lies on the theta divisor of $B^1$.

For clarity, the ordinarity implication uses the exact
sequence
\[
 0\longrightarrow L\longrightarrow F_*F^*L
   \longrightarrow B^1\otimes L\longrightarrow0.
\]
Since $H^0(L)=H^0(L^2)=0$, the just-proved vanishing
makes Frobenius an isomorphism between the two
$g-1$ dimensional twisted $H^1(\mathcal O)$ spaces.
These are precisely the nontrivial character summands
of $H^1(D,\mathcal O_D)$, as
$\pi_*\mathcal O_D=\mathcal O_C\oplus L\oplus L^2$.
Their stable Frobenius rank is $2(g-1)$, which proves
the assertion. This is also the rank-one instance of
Raynaud's ordinarity criterion in \S1.
\citep[\S1, before Proposition~1]{LA_Raynaud2000}

Here are the details of the reflection assertion in
characteristic two. Let $\delta$ generate the cyclic
deck group, and let $\iota$ be the hyperelliptic
involution on $C$. Since $\iota^*L=L^{-1}$,
$\iota$ lifts to an automorphism $\tau$ of $D$ with
$\tau\delta\tau^{-1}=\delta^{-1}$.
Write $\tau^2=\delta^a$. The element $\tau$ commutes
with its square and also sends $\delta^a$ to
$\delta^{-a}$, so $a=-a$ in $\Z/3$ and $a=0$.
Thus $\tau$ is an involution. It fixes exactly one
of the three points over $P$, since it acts there
as a reflection of the three-element deck torsor.
It has no other fixed points, as $\iota$ has none.
The \'etale map $D\to C$ identifies the completed
local rings at that fixed point and at $P$.
The local different for $D\to H:=D/\langle\tau\rangle$
therefore has exponent $2g+2$, the same as that of
$C\to\mathbb P^1$. Consequently
\[
 6g-6=2g_D-2=2(2g_H-2)+(2g+2),
 \qquad g_H=g-1.
\]

Let $A,A'$ be the pullback images in $J_D$ of the
Jacobians of the quotients by $\tau$ and by a distinct
conjugate reflection. Both lie in the Prym: their
norm to $J_C$ is fixed by the induced hyperelliptic
action $[-1]$, hence has finite image and is zero
on these connected abelian varieties. Their
intersection is finite. Indeed an element fixed by
the two reflections is fixed by $\delta$ as well.
Its norm $z$ to $J_C$ satisfies $2z=0$, while the
pullback of $z$ equals $3$ times that element; hence
the intersection is killed by six. Pullback from
each quotient has finite kernel by its pullback--norm
identity. Both images thus have dimension $g-1$,
and their addition gives an isogeny
\[
 J_H\times J_H\longrightarrow\Pr(D/C).
\]
For $g=3$ the quotient therefore has genus two and
is ordinary.
The inequality used above fails for $g\ge4$; no
assertion in those genera follows from this proof.
\end{proof}

\begin{claim}{B.LIT9}{Proved}{A fixed simply branched function cannot supply the missing fibres}
Let $k$ be algebraically closed, let $f:C\to\mathbb P^1_k$
be finite separable, and assume that for every
$t\in\mathbb P^1(k)$,
\[
 \gcd\{e_P(f):P\in f^{-1}(t)\}=1.
\]
Let $C'\to C$ be any finite \'etale cover and
$h:\mathbb P^1\to\mathbb P^1$ any finite separable map.
The composite $h\circ f\circ\pi$ cannot have three
distinct fibres in each of which all ramification
indices are divisible by three. It cannot have four
distinct fibres in each of which they are all even.
More generally, uniform fibre divisibilities $m_i$
at distinct values require
\[
 \sum_i(1-1/m_i)<2.
\]
\end{claim}

\begin{proof}
The \'etale pullback preserves every ramification index
over the algebraically closed ground field. For a
fixed $t$ over $u$ the gcd of the indices of the
composite at points over $t$ is therefore $e_t(h)$.
Taking the gcd over all $t\in h^{-1}(u)$ shows that
any uniform divisor $m_i$ of the composite fibre
already divides every $e_t(h)$ in that fibre of $h$.
Write $n=\deg h$. There are at most $n/m_i$ geometric
points in such a fibre. Their contribution to the
different is at least
$\sum_{t\mid u_i}(e_t(h)-1)\ge n(1-1/m_i)$,
also in the presence of wild ramification.
Riemann--Hurwitz for $h:\mathbb P^1\to\mathbb P^1$
now gives
\[
 n\sum_i(1-1/m_i)\le2n-2<2n.
\]
The claimed patterns each give a sum of two and are
impossible. In particular, the hypothesis holds for
a simply branched map of degree at least three:
every branch fibre has an unramified point. Thus
rational postcomposition and \'etale pullback of that
fixed function cannot complete the general B5 route;
a new function on the covering curve would be needed.
\end{proof}

\begin{claim}{B.LIT10}{Proved}{Three uniformly divisible fibres with a supersingular Fermat quotient}
Let $C/\overline{\F}_p$ be a smooth proper curve admitting
a finite separable map to $\mathbb P^1$ with three
distinct fibres in which every ramification index is
divisible by the same odd integer $r\ge3$.
Suppose $r\mid p^a+1$ for some $a\ge1$.
Then $C$ has a connected cyclic \'etale cover of degree
dividing $r$ which dominates a supersingular curve of
genus $(r-1)/2$. Its geometric Galois closure therefore
has order prime to $p$.
In characteristic two this applies to every cyclic
prime-degree cover of $\mathbb P^1$ of genus at least
two whose degree $\ell$ satisfies that
$\operatorname{ord}_\ell(2)$ is even.
\end{claim}

\begin{proof}
Send the three specified values to $0,1,\infty$.
The smooth proper model
\[
 H_r:\quad y^r=x(x-1)
\]
is a cyclic cover of degree $r$, totally ramified
at those three values and nowhere else. Since $r$
is odd, its genus is $(r-1)/2$. Normalizing
$C\times_{\mathbb P^1}H_r$ and taking any connected
component gives a cyclic cover $D\to C$ of degree
dividing $r$. At each of the three relevant fibres,
the tame local ramification index $r$ divides that
of $C\to\mathbb P^1$, so the normalized pullback is
unramified. It is unramified elsewhere as well.
Thus $D\to C$ is \'etale, and $D\to H_r$ is dominant.

Put $q=p^a$. The Fermat curve of degree $r$ maps
to $H_r$ by
\[
 U^r+V^r=1,\qquad x=U^r,\quad y=-UV.
\]
It is in turn dominated by the Hermitian Fermat
curve $U^{q+1}+V^{q+1}=1$, by taking coordinate
powers $(q+1)/r$. This latter curve has genus
$q(q-1)/2$ and exactly $q^3+1$ points over
$\F_{q^2}$. Indeed the norm map to $\F_q$ has
one preimage of zero and $q+1$ preimages of each
nonzero value. Counting affine solutions gives
$2(q+1)+(q-2)(q+1)^2=q^3-q$, and there are
$q+1$ points at infinity. This attains the upper
Weil bound, so every Frobenius root equals $-q$.
The curve is therefore supersingular. Pullback
and norm inject the Jacobian of every dominated
curve in the isogeny category; hence $H_r$ is
supersingular too.

For the last assertion, a cyclic cover of odd prime
degree $\ell$ has all ramification indices equal
to $\ell$ at its branch values. Genus at least two
forces at least three such values. An even order
of $2$ modulo $\ell$ gives
$2^{\operatorname{ord}_\ell(2)/2}=-1$ modulo
$\ell$, so the proved criterion applies. Degrees
$5,11,13,17,19$ give examples beyond the cubic
case. Degree $7$ has order three and is outside
this criterion; no inference for arbitrary
characteristic-two curves follows.
\end{proof}
